\documentclass{ximera}
\usepackage{hyperref}
\hypersetup{colorlinks=true, linkcolor=blue, filecolor=magenta, urlcolor=cyan,}
\urlstyle{same}


\begin{document}
\section*{MATH 1075: Precollege Mathematics II}
\section*{Spring 2026 Course Syllabus}
\section*{Course material}
\section*{Textbook + ALEKS subscription}
Beginning and Intermediate Algebra by Miller, O'Neill, and Hyde (6th edition)

★ All students in the course are required to have access to the ALEKS homework system that is used with the course. The online homework/ study system includes a (searchable) electronic version of the textbook.

★ The textbook and/or courseware for this course is being provided via CarmenBooks. Through CarmenBooks, students obtain publisher materials electronically through CarmenCanvas, saving them up to $80 \%$ per title. To avail the discounted price, students should access the textbook system via Canvas only.

\section*{Calculators}
Students are required to have a scientific or graphing calculator; a graphing calculator is strongly recommended (TI-83 or TI-84 would be best). Calculators with a Computer Algebra System (e.g., TI-89 and TI-92) are prohibited from use on quizzes and exams. If you are not sure about your calculator, ask your lecturer or recitation instructor as soon as possible.

\section*{Course webpage, E-mail \& Canvas}
\begin{itemize}
  \item The mathematics department Math 1075 course webpage can be accessed at \href{http://math.osu.edu/courses/1075}{http://math.osu.edu/courses/1075}.
  \item You are responsible for information contained in e-mail messages sent to your official OSU e-mail address: yourlastname.\#@osu.edu. You should check your e-mail at least once per day to avoid missing time sensitive information.
  \item Canvas is a web-based course tool that allows you to access all course material. Course communications will primarily arrive via Canvas messages and announcements. You must enable notifications from Canvas on your preferred device to participate in the course.
\end{itemize}

\section*{Grading Scale (Percent)}
Course grades will be allotted based on the following grading scale. This course is not curved, and no further adjustments may be made at the end of the semester.

\begin{center}
\begin{tabular}{c|c|c|c|c|c|c|c|c|c}
A & A- & B+ & B & B- & C+ & C & C- & D+ & D \\
\hline
90 & 87 & 83 & 80 & 77 & 73 & 70 & 67 & 63 & 60 \\
\hline
\end{tabular}
\end{center}

\section*{Course Grade}
The course grade will be computed out of 100 points, comprised of

\begin{center}
\begin{tabular}{l|l}
ALEKS Homework & 15 points \\
Quizzes (in recitation) & 15 points \\
Midterm 1 & 15 points \\
Midterm 2 & 15 points \\
Midterm 3 & 15 points \\
Final Exam & 25 points \\
Extra Credit & 2 points \\
\end{tabular}
\end{center}

\begin{itemize}
  \item HOMEWORK is due online via ALEKS. See the course calendar for deadlines of each ALEKS homework assignment. Students are responsible for completing ALEKS homework by 11:59pm on the due date. The ALEKS homework system does not accept late homework.
  \item QUIZZES will be administered in-recitation . The date of each quiz as well as sections to study for each quiz are given in the course calendar. There are ten (10) quizzes throughout this course, and your quiz average is calculated with your eight ( 8 ) highest quiz scores. Makeup quizzes will only be provided for university sanctioned conflicts, and missed quizzes will automatically count as one of the dropped quizzes when appropriate.
  \item ATTENDANCE will be taken in every recitation that is not an exam day or a quiz day. This will count for up to 2 points of extra credit toward your final grade. The extra credit will be applied at the end of the semester.
\end{itemize}

\section*{Examinations}
There will be three midterm exams and a cumulative and comprehensive final exam. The locations of examinations will be made available a week before each examination. All exams will be in-person.

Midterm 1 . . . Tuesday, Feb 32026 . . . 5:20pm-6:15pm\\
Midterm 2 . . . Tuesday, Mar 32026 . . . 5:20pm-6:15pm\\
Midterm 3 . . . Tuesday, Mar 312026 . . . 5:20pm-6:15pm\\
Final Exam ... Friday, May 12026 ... 6:00pm-7:45pm

\section*{EARLY EXAMS ARE NOT PERMITTED UNDER ANY CIRCUMSTANCES in MATH 1075.}
Makeup examinations will only be offered to students with UNIVERSITY SANCTIONED CONFLICTS.\\
Makeup exams will take place the following mornings at 8am and will require the permission of your lecturer before the scheduled examinations take place.

\section*{Help with the course}
★ Your lecturer and recitation instructors will hold regular office hours each week for individual help.

★ The Math Stat Learning Center (MSLC) offers free drop-in and appointment tutoring. Everyone can benefit from tutoring. The MSLC's drop-in tutor rooms are a great place to work on math homework or study for exams. Students often use the space like a library with the added benefit of a tutor or peers nearby. MSLC tutors focus not only on helping you solve the problem at hand, but also work with you to build your understanding and knowledge to prepare you for exams. (\href{https://mslc.osu.edu}{https://mslc.osu.edu})

Math 1075 drop-in tutoring is held in

\begin{itemize}
  \item CH004 (Cockins Hall) from 10:20am 7:30pm, Monday - Thursday
  \item MA010 (Math Building) from 1pm-4pm on Fridays and Sundays.
\end{itemize}

Appointment based tutoring is also available at the MSLC.

\section*{Disability services statement}
Students with disabilities that have been certified by Student Life Disabilities Services (SLDS) will be appropriately accommodated and should inform the instructor as soon as possible via the AIM portal. Contact information of SLDS is as follows.

\begin{itemize}
  \item Website: \href{http://slds.osu.edu/}{http://slds.osu.edu/}
  \item Email: \href{mailto:slds@osu.edu}{slds@osu.edu}
  \item Phone: 614-292-3307
  \item VRS: 614-500-4445
  \item Address: 98 Baker Hall, 113 W. 12th Avenue, Columbus, OH 43210
\end{itemize}

\section*{Catalog description}
Algebraic, rational, and radical expressions; functions and graphs; quadratic equations; absolute value; inequalities; and applications.

\section*{Prerequisites}
Math Placement Level S, a grade of C- or better in Math 1050, or credit for Math 75 or 1074. Not open to students with credit for any higher numbered math class, or for any quarter math class numbered higher than 75 .

\section*{GE Information}
This mathematics course can be used, depending on your degree program, to satisfy the Quantitative Reasoning: Basic Computation category of the General Education Requirement (GE). The goals and learning objectives for this category are:

\begin{itemize}
  \item Goals: Students develop skills in quantitative literacy and logical reasoning, including the ability to identify valid arguments, and use mathematical models.
  \item Learning objectives: Students demonstrate computational skills and familiarity with algebra and geometry and apply these skills to practical problems.
\end{itemize}

\section*{Academic misconduct statement}
It is the responsibility of the Committee on Academic Misconduct to investigate or establish procedures for the investigation of all reported cases of student academic misconduct. The term academic misconduct includes all forms of student academic misconduct wherever committed; illustrated by, but not limited to, cases of plagiarism and dishonest practices in connection with examinations. Instructors shall report all instances of alleged academic misconduct to the committee. (Faculty Rule 3335-5-48.7). For additional information, see the Code of Student Conduct at \href{http://studentlife.osu.edu/csc/}{http://studentlife.osu.edu/csc/}

\section*{Religious Accomodations}
Ohio State has had a longstanding practice of making reasonable academic accommodations for students' religious beliefs and practices in accordance with applicable law. In 2023, Ohio State updated its practice to align with new state legislation. Under this new provision, students must be in early communication with their instructors regarding any known accommodation requests for religious beliefs and practices, providing notice of specific dates for which they request alternative accommodations within 14 days after the first instructional day of the course. Instructors in turn shall not question the sincerity of a student's religious or spiritual belief system in reviewing such requests and shall keep requests for accommodations confidential.\\
With sufficient notice, instructors will provide students with reasonable alternative accommodations with regard to examinations and other academic requirements with respect to students' sincerely held religious beliefs and practices by allowing up to three absences each semester for the student to attend or participate in religious activities. Examples of religious accommodations can include, but are not limited to, rescheduling an exam, altering the time of a student's presentation, allowing make-up assignments to substitute for missed class work, or flexibility in due dates or research responsibilities. If concerns arise about a requested accommodation, instructors are to consult their tenure initiating unit head for assistance.\\
A student's request for time off shall be provided if the student's sincerely held religious belief or practice severely affects the student's ability to take an exam or meet an academic requirement and the student has notified their instructor, in writing during the first 14 days after the course begins, of the date of each absence. Although students are required to provide notice within the first 14 days after a course begins, instructors are strongly encouraged to work with the student to provide a reasonable accommodation if a request is made outside the notice period. A student may not be penalized for an absence approved under this policy. If students have questions or disputes related to academic accommodations, they should contact their course instructor, and then their department or college office. For questions or to report discrimination or harassment based on religion, individuals should contact the Civil Rights Compliance Office

\section*{Intellectual Diversity}
Ohio State is committed to fostering a culture of open inquiry and intellectual diversity within the classroom. This course will cover a range of information and may include discussions or debates about controversial issues, beliefs, or policies. Any such discussions and debates are intended to support understanding of the approved curriculum and relevant course objectives rather than promote any specific point of view. Students will be assessed on principles applicable to the field of study and the content covered in the course. Preparing students for citizenship includes helping them develop critical thinking skills that will allow them to reach their own conclusions regarding complex or controversial matters.


\end{document}